\subsection{Brion's integration map}

Brion in \cite{Brion} defined the isomorphism 
\[  H^n(\Delta) \to K \]
in terms of piecewise polynomial functions on the fan $\Delta$. We describe this map in the more general case where the field $K$ is not necessarily  $\RR$, and $\Delta$ is a pseudo-manifold.

The isomorphism depends on a fixed volume form 
\[ t_1\wedge t_2 \wedge \cdots \wedge t_n \in \Lambda^n V^*,\]
and an orientation on $\Delta$ over the field $K$.  

Let $\sigma$ be a maximal simplex in $\Delta$, and let $v_{j_1},\ldots,v_{j_n}$ be an ordering of its vertices given by the orientation. (If the characteristic of $K$ is $2$, then any ordering is allowed.)
Using the isomorphism~\eqref{eq-isom}, define the polynomial $\chi_\sigma \in K[t_1,\ldots, t_n]$ as
\[ \chi_\sigma = c_\sigma x_{j_1} x_{j_2} \cdots x_{j_n},\]
where the constant $c_\sigma \in K$ is such that 
\[ c_\sigma x_{j_1} \wedge x_{j_2} \wedge \cdots \wedge x_{j_n} = t_1\wedge t_2 \wedge \cdots \wedge t_n.\]
One can compute that 
 \begin{equation} \label{eq-dets}
  c_\sigma= \det\sigma = \det \begin{bmatrix} 
 \a{1, j_1} & \a{1, j_2} & \ldots & \a{1, j_n} \\
 \vdots & \vdots & \ddots & \vdots \\
  \a{n, j_1} & \a{n, j_2} & \ldots & \a{n, j_n}
  \end{bmatrix}.
  \end{equation}

The following lemma was proved in \cite[Theorem~2.2]{Brion} in the case where the field is $\RR$ and $\Delta$ is a complete fan. The proof for an arbitrary field $K$ and an oriented pseudo-manifold $\Delta$ is the same, so we recall it.  

\begin{lemma} 
Let $\Delta$ be an oriented pseudo-manifold of dimension $n-1$ over $k$. Consider
\begin{align} \label{eq-int}
 \pi_{\Delta}: \cA(\Delta) & \to K(t_1, \ldots, t_n) \\ \nonumber
 f & \mapsto  \sum_{\sigma \in \Delta_{n-1}} \frac{f_\sigma}{\chi_\sigma}.
 \end{align} 
Then the image of $\pi_{\Delta}$ lies in $K[t_1, \ldots, t_n]$ and the induced map
\[ \pi_{\Delta}: \cA(\Delta)  \to K[t_1, \ldots, t_n]\]
is a degree $-n$ homomorphism of graded $K[t_1,\ldots, t_n]$ modules. The map $\pi_{\Delta}$ in degree $n$ defines an isomorphism
\begin{equation*} 
 \pi_{\Delta}: H^n(\Delta) \to K.
 \end{equation*}
\end{lemma}

\begin{proof}
We start by proving the last statement of the lemma. Recall that $t_i$ act on $\cA(\Delta)$ by multiplication with $\theta_i$. It follows that $\pi_{\Delta}$ maps $(\theta_1,\ldots,\theta_n) \cA^{n-1}(\Delta)$ to zero because it maps $\cA^{n-1}(\Delta)$ to elements of degree $-1$ in $K[t_1,\ldots, t_n]$. Hence $\pi_{\Delta}$ in degree $n$ factors through $H^n(\Delta)$. This map is nonzero because the piecewise polynomial function $\{ f_\sigma\}$ such that $f_{\sigma_0} = \chi_{\sigma_0}$ and $f_\sigma=0$ for $\sigma\neq \sigma_0$ for some fixed $\sigma_0$ maps to $1\in K$.

The map $\pi_{\Delta}$ is clearly a homomorphism of $K[t_1,\ldots, t_n]$-modules: when we multiply $f= \{f_\sigma \}$ with $\theta_i$, we multiply each $f_\sigma(t_1,\ldots,t_n)$ with $t_i$. The map decreases degree by $n$ because all $\chi_\sigma$ are homogeneous polynomials of degree $n$. It remains to show that the image of $\pi_{\Delta}$ lies in the polynomial ring.

Each $\chi_\sigma$ is a product of linear functions that vanish on the $n$ hyperplanes in $V$ spanned by the facets of $\sigma$. This implies that the rational function $f_\sigma/\chi_\sigma$ can have at worst simple poles along these hyperplanes. Fix one $(n-2)$-dimensional simplex $\tau$ and let $H$ be the hyperplane it spans. Let $\sigma_1$ and $\sigma_2$ be the two $(n-1)$-dimensional simplices containing $\tau$. Now it suffices to prove that the residues of  $f_{\sigma_1}/\chi_{\sigma_1}$ and $f_{\sigma_2}/\chi_{\sigma_2}$ along the hyperplane $H$ sum to zero. This implies that all poles cancel and the image of $\pi_{\Delta}$ is a polynomial.

Consider $\chi_{\sigma_1} =  (\det \sigma_1) x_{j_1} x_{j_2} \cdots x_{j_n}$, where we take $t=(\det \sigma_1) x_{j_1}$ as the parameter that vanishes on $H$. The residue of $f_{\sigma_1}/\chi_{\sigma_1}$ with respect to the parameter $t$ is then 
\[ \frac{f_{\sigma_1}}{x_{j_2} \cdots x_{j_n}} \pmod t.\]
Working mod $t$ means that we restrict the rational function to the hyperplane $H$.

Now consider $\chi_{\sigma_2} = (\det \sigma_2) x_{j_1'} x'_{j_2} \cdots x'_{j_n}$, where $(\det \sigma_2) x_{j_1'}$ vanishes on $H$ and $x'_{j_2}, \ldots, x'_{j_n}$ are equal to $x_{j_2}, \ldots, x_{j_n}$ when restricted to the hyperplane $H$. From the normalization condition
\begin{align*}
  (\det \sigma_1) x_{j_1}\wedge x_{j_2} \wedge \cdots \wedge x_{j_n} &= -(\det \sigma_2) x_{j_1'}\wedge x'_{j_2} \wedge \cdots \wedge x'_{j_n} \\
    &= -(\det \sigma_2) x_{j_1'}\wedge x_{j_2} \wedge \cdots \wedge x_{j_n},
\end{align*}
it follows that 
\[ t= (\det \sigma_1) x_{j_1} = - (\det \sigma_2) x_{j_1'}.\]
The residue of $f_{\sigma_2}/\chi_{\sigma_2}$ with respect to the parameter $t$ is
\[ -\frac{f_{\sigma_2}}{x'_{j_2} \cdots x'_{j_n}} \equiv - \frac{f_{\sigma_2}}{x_{j_2} \cdots x_{j_n}} \pmod t.\]
This is equal to the negative of the residue of $f_{\sigma_1}/\chi_{\sigma_1}$ because $f_{\sigma_1}$ and $f_{\sigma_2}$ restrict to the same polynomial on $H$.
\end{proof}
